\chapter{Jednostajna ciągłość}
\label{ma1:ch:uniform-continuity}

W definicji ciągłości promień $\delta$ może zależeć od punktu dziedziny. Jednostajna ciągłość wymaga, aby przy ustalonej dokładności obrazów jeden promień działał dla wszystkich par dostatecznie bliskich argumentów. Pojęcie to wiąże ciągłość z metryką, zachowaniem ciągów Cauchy'ego oraz przedłużaniem odwzorowań.

\section{Definicja i kryterium ciągowe}
\begin{definicja}
\label{ma1:def:uniform-continuity}
Niech $(X,d_X)$ i $(Y,d_Y)$ będą przestrzeniami metrycznymi, $A\subseteq X$ i $f:A\to Y$. Odwzorowanie $f$ jest jednostajnie ciągłe na $A$, jeżeli
\[
 \forall\varepsilon>0\ \exists\delta>0\ \forall x,y\in A:
 \quad d_X(x,y)<\delta\ \Longrightarrow\ d_Y(f(x),f(y))<\varepsilon.
\]
\end{definicja}
Dla porównania ciągłość na $A$ ma postać
$\forall x\in A\ \forall\varepsilon>0\ \exists\delta>0\ \forall y\in A:
d_X(x,y)<\delta\Rightarrow d_Y(f(x),f(y))<\varepsilon$.
Różnica polega na położeniu kwantyfikatora dotyczącego $x$. Każde odwzorowanie jednostajnie ciągłe jest ciągłe: przy ustalonym $x$ wybieramy promień z definicji jednostajnej ciągłości.

\begin{twierdzenie}
\label{ma1:thm:uniform-sequential}
Niech $X,Y$ będą przestrzeniami metrycznymi, $A\subseteq X$ i $f:A\to Y$. Odwzorowanie $f$ jest jednostajnie ciągłe wtedy i tylko wtedy, gdy dla każdych dwóch ciągów $(x_n),(y_n)$ w $A$ z warunkiem $d_X(x_n,y_n)\to0$ zachodzi $d_Y(f(x_n),f(y_n))\to0$.
\end{twierdzenie}
\begin{proof}
Jeżeli $f$ jest jednostajnie ciągłe, to dla $\varepsilon>0$ wybieramy $\delta>0$ z definicji. Od pewnego indeksu $d_X(x_n,y_n)<\delta$, więc $d_Y(f(x_n),f(y_n))<\varepsilon$.

Odwrotnie, jeśli $f$ nie jest jednostajnie ciągłe, zaprzeczenie definicji daje $\varepsilon_0>0$ takie, że dla każdego $\delta>0$ istnieją $x,y\in A$ spełniające $d_X(x,y)<\delta$ i $d_Y(f(x),f(y))\ge\varepsilon_0$. Dla $\delta=1/n$, $n\ge1$, wybieramy parę $(x_n,y_n)$. Odległości argumentów dążą do zera, ale odległości obrazów nie dążą do zera. Otrzymujemy sprzeczność.
\end{proof}
W tym kryterium nie zakładamy, że którykolwiek z dwóch ciągów argumentów jest zbieżny.

\begin{przyklad}
\label{ma1:ex:square-not-uniform}
Funkcja $f(x)=x^2$ jest ciągła na $\R$, ale nie jest jednostajnie ciągła. Dla $x_n=n$, $y_n=n+1/n$ mamy $|y_n-x_n|=1/n\to0$, natomiast
$|f(y_n)-f(x_n)|=2+1/n^2\to2$.
Na każdym przedziale $[-R,R]$, $R>0$, jest jednak jednostajnie ciągła, gdyż
$|x^2-y^2|\le2R|x-y|$, więc wystarczy wybrać $\delta=\varepsilon/(2R)$.
\end{przyklad}
\begin{przyklad}
\label{ma1:ex:reciprocal-not-uniform}
Funkcja $f:(0,1]\to\R$, $f(x)=1/x$, jest ciągła, lecz dla $x_n=1/n$, $y_n=1/(n+1)$ odległość argumentów wynosi $1/(n(n+1))\to0$, podczas gdy $|f(x_n)-f(y_n)|=1$. Zatem nie jest jednostajnie ciągła nawet na ograniczonej dziedzinie.
\end{przyklad}

\section{Twierdzenie Heinego--Cantora}
\begin{twierdzenie}[Heinego--Cantora]
\label{ma1:thm:heine-cantor}
Niech $X,Y$ będą przestrzeniami metrycznymi, $K\subseteq X$ będzie zbiorem ciągowo zwartym i $f:K\to Y$ odwzorowaniem ciągłym. Wówczas $f$ jest jednostajnie ciągłe.
\end{twierdzenie}
\begin{proof}
Przypuśćmy, że $f$ nie jest jednostajnie ciągłe. Z zaprzeczenia definicji istnieją $\varepsilon_0>0$ i ciągi $(x_n),(y_n)$ w $K$ takie, że $d_X(x_n,y_n)<1/n$, ale $d_Y(f(x_n),f(y_n))\ge\varepsilon_0$.

Ze zwartości ciągowej wybieramy podciąg $x_{n_k}\to p\in K$. Z nierówności trójkąta
$d_X(y_{n_k},p)\le d_X(y_{n_k},x_{n_k})+d_X(x_{n_k},p)\to0$.
Oba podciągi dążą do tego samego punktu $p$ w dziedzinie. Ciągłość $f$ w $p$ daje $f(x_{n_k})\to f(p)$ i $f(y_{n_k})\to f(p)$. Wobec tego
\[
 d_Y(f(x_{n_k}),f(y_{n_k}))
 \le d_Y(f(x_{n_k}),f(p))+d_Y(f(p),f(y_{n_k}))\to0,
\]
co przeczy dolnemu ograniczeniu $\varepsilon_0$.
\end{proof}
\begin{wniosek}
\label{ma1:cor:interval-uniform}
Jeżeli $a,b\in\R$, $a\le b$, i $f:[a,b]\to Y$ jest ciągłe, gdzie $Y$ jest przestrzenią metryczną, to $f$ jest jednostajnie ciągłe.
\end{wniosek}
\begin{proof}
Przedział $[a,b]$ jest ciągowo zwarty na mocy lematu \ref{ma1:lem:interval-seqcompact}. Stosujemy twierdzenie \ref{ma1:thm:heine-cantor}.
\end{proof}
Ograniczoność dziedziny nie zastępuje zwartości ciągowej (przykład \ref{ma1:ex:reciprocal-not-uniform}), podobnie jak sama domkniętość (przykład \ref{ma1:ex:square-not-uniform}). Z drugiej strony zwartość ciągowa nie jest konieczna do jednostajnej ciągłości konkretnego odwzorowania, np. identyczność $\R\to\R$ jest jednostajnie ciągła.

\section{Warunki Lipschitza i H\"oldera}
\begin{definicja}
\label{ma1:def:holder-lipschitz}
Niech $X,Y$ będą przestrzeniami metrycznymi, $A\subseteq X$ i $f:A\to Y$. Dla $0<\alpha\le1$ mówimy, że $f$ spełnia warunek H\"oldera z wykładnikiem $\alpha$, jeżeli istnieje $C\ge0$ takie, że
\[
 d_Y(f(x),f(y))\le C d_X(x,y)^\alpha
 \qquad(x,y\in A).
\]
Przypadek $\alpha=1$ nazywamy warunkiem Lipschitza, a stałą $C$ stałą Lipschitza.
\end{definicja}
\begin{twierdzenie}
\label{ma1:thm:holder-uniform}
Każde odwzorowanie spełniające warunek H\"oldera z wykładnikiem $0<\alpha\le1$ jest jednostajnie ciągłe.
\end{twierdzenie}
\begin{proof}
Jeśli $C=0$, odległość obrazów dowolnych dwóch punktów wynosi zero; dowolny $\delta>0$ jest wystarczający. Jeżeli $C>0$, dla $\varepsilon>0$ wybieramy $\delta=(\varepsilon/C)^{1/\alpha}$. Z $d_X(x,y)<\delta$ wynika $C d_X(x,y)^\alpha<\varepsilon$.
\end{proof}
\begin{przyklad}
\label{ma1:ex:sqrt-holder}
Dla $0\le y\le x$ mamy $\sqrt{xy}\ge y$, więc
$(\sqrt x-\sqrt y)^2=x+y-2\sqrt{xy}\le x-y$.
Po zamianie $x,y$, jeśli jest konieczna, otrzymujemy dla wszystkich $x,y\ge0$
$|\sqrt x-\sqrt y|\le\sqrt{|x-y|}$.
Funkcja pierwiastkowa na $[0,+\infty)$ spełnia warunek H\"oldera z wykładnikiem $1/2$ i stałą $1$. Nie spełnia warunku Lipschitza, ponieważ dla $x=1/n^2$ i $y=0$ iloraz $|\sqrt x-\sqrt y|/|x-y|=n$ jest nieograniczony.
\end{przyklad}
\begin{przyklad}
\label{ma1:ex:uniform-no-holder}
Niech $A=\{0\}\cup\{2^{-n^2}:n\ge1\}$ i niech $f(0)=0$, $f(2^{-n^2})=1/n$. Zbiór $A$ jest ciągowo zwarty: jeżeli pewna wartość występuje nieskończenie wiele razy w ciągu w $A$, istnieje podciąg stały; w przeciwnym przypadku można wybrać podciąg wyrazów $2^{-n_k^2}$ o $n_k\to+\infty$, zbieżny do $0$. Każdy dodatni punkt $A$ jest izolowany, ponieważ kolejne dodatnie elementy tworzą ściśle malejący ciąg zbieżny do zera. W zerze $f$ jest ciągła: dla $\varepsilon>0$ wybieramy $N$ z $1/N<\varepsilon$, a następnie $\delta=2^{-N^2}$; jeśli $x\in A$ i $0<x<\delta$, to $x=2^{-n^2}$ dla $n>N$, więc $|f(x)|<\varepsilon$.
Z twierdzenia Heinego--Cantora $f$ jest jednostajnie ciągła. Dla dowolnego $\alpha>0$ iloraz
$|f(2^{-n^2})-f(0)|/|2^{-n^2}|^\alpha=2^{\alpha n^2}/n$ dąży do $+\infty$ (od pewnego indeksu $\alpha n\ge1$, więc jest nie mniejszy niż $2^n/n\to+\infty$). Nie istnieje zatem stała H\"oldera dla żadnego $\alpha\in(0,1]$.
\end{przyklad}

\begin{center}
\begin{tikzpicture}[node distance=6mm,every node/.style={font=\small,align=center},>=Stealth]
\node[draw,rounded corners] (l) {Warunek Lipschitza};
\node[draw,rounded corners,below=of l] (h) {Warunek H\"oldera dla pewnego $\alpha\in(0,1]$};
\node[draw,rounded corners,below=of h] (u) {Jednostajna ciągłość};
\node[draw,rounded corners,below=of u] (c) {Ciągłość};
\draw[->] (l)--(h);
\draw[->] (h)--(u);
\draw[->] (u)--(c);
\end{tikzpicture}

\small Rysunek: implikacje między warunkami regularności; żadna odwrotność nie zachodzi ogólnie.
\end{center}

\section{Moduł ciągłości}
\begin{definicja}
\label{ma1:def:modulus-continuity}
Niech $X,Y$ będą przestrzeniami metrycznymi, $A\subseteq X$ będzie niepusty i $f:A\to Y$. Modułem ciągłości odwzorowania $f$ nazywamy funkcję $\omega_f:[0,+\infty)\to[0,+\infty]$ określoną przez
\[
 \omega_f(t)=\sup\{d_Y(f(x),f(y)):x,y\in A,\ d_X(x,y)\le t\}.
\]
Jeżeli zbiór odległości jest nieograniczony, przyjmujemy supremum równe $+\infty$.
\end{definicja}
Zbiór w definicji zawiera $0$, ponieważ możemy wybrać $x=y$. W szczególności $\omega_f(0)=0$. Dla $s\le t$ zbiór par odpowiadający $s$ zawiera się w zbiorze par odpowiadającym $t$, więc $\omega_f(s)\le\omega_f(t)$.
\begin{twierdzenie}
\label{ma1:thm:modulus-uniform}
Niech $f:A\to Y$, gdzie $A$ jest niepustym podzbiorem przestrzeni metrycznej $X$, a $Y$ jest przestrzenią metryczną. Odwzorowanie $f$ jest jednostajnie ciągłe wtedy i tylko wtedy, gdy $\lim_{t\to0^+}\omega_f(t)=0$, tzn. dla każdego $\varepsilon>0$ istnieje $\eta>0$ takie, że $\omega_f(t)<\varepsilon$ dla $0<t<\eta$.
\end{twierdzenie}
\begin{proof}
Jeżeli $f$ jest jednostajnie ciągłe, dla $\varepsilon/2$ wybieramy $\delta>0$. Dla $0<t<\delta$ wszystkie pary z $d_X(x,y)\le t$ spełniają $d_Y(f(x),f(y))<\varepsilon/2$. Ich supremum jest nie większe niż $\varepsilon/2<\varepsilon$.
Odwrotnie, wybierzmy $\eta>0$ dla $\varepsilon$ z warunku granicznego. Gdy $0<d_X(x,y)<\eta$, przyjmujemy $t=d_X(x,y)$ i otrzymujemy $d_Y(f(x),f(y))\le\omega_f(t)<\varepsilon$. Gdy $d_X(x,y)=0$, mamy $x=y$ i odległość obrazów jest zerowa.
\end{proof}
\begin{przyklad}
\label{ma1:ex:modulus-exact}
Dla funkcji $f(x)=\sqrt x$ na $[0,+\infty)$ mamy $\omega_f(t)=\sqrt t$ dla $t\ge0$. Górne oszacowanie wynika z przykładu \ref{ma1:ex:sqrt-holder}, a para $(0,t)$ daje równość. Dla identyczności na $\R$ analogicznie $\omega_f(t)=t$.
Moduł nie musi być skończony dla dużych $t$ nawet przy jednostajnej ciągłości: na $\N$ z metryką dyskretną funkcja $f(n)=n$ jest jednostajnie ciągła (wystarczy $\delta=1$), lecz $\omega_f(t)=+\infty$ dla $t\ge1$.
\end{przyklad}

\section{Działania na odwzorowaniach jednostajnie ciągłych}
\begin{twierdzenie}
\label{ma1:thm:uniform-operations}
Niech $X,Y,Z$ będą przestrzeniami metrycznymi, $A\subseteq X$, $B\subseteq Y$.
\begin{enumerate}[label=(\roman*)]
\item Ograniczenie jednostajnie ciągłego odwzorowania $f:A\to Y$ do dowolnego podzbioru $A$ jest jednostajnie ciągłe.
\item Jeżeli $f:A\to B$ i $g:B\to Z$ są jednostajnie ciągłe, to $g\circ f$ jest jednostajnie ciągłe.
\item Jeżeli $f,g:A\to K$, gdzie $K=\R$ lub $K=\C$, są jednostajnie ciągłe, to $\lambda f+\mu g$ jest jednostajnie ciągłe dla $\lambda,\mu\in K$. Jeżeli ponadto $f,g$ są ograniczone, to $fg$ jest jednostajnie ciągłe.
\item Jeżeli $g:A\to K$ jest jednostajnie ciągłe i istnieje $m>0$ takie, że $|g(x)|\ge m$ dla wszystkich $x\in A$, to $1/g$ jest jednostajnie ciągłe.
\end{enumerate}
\end{twierdzenie}
\begin{proof}
(i) Ten sam promień działa na podzbiorze. (ii) Dla $\varepsilon>0$ wybieramy $\rho>0$ z jednostajnej ciągłości $g$, a dla $\rho$ wybieramy $\delta>0$ z jednostajnej ciągłości $f$. Wtedy $d_X(x,y)<\delta$ implikuje $d_Y(f(x),f(y))<\rho$, a następnie $d_Z(g(f(x)),g(f(y)))<\varepsilon$.

(iii) Wybierzmy wspólny promień (minimum dwóch promieni), dla którego różnice wartości $f,g$ są mniejsze niż $\varepsilon/(1+|\lambda|+|\mu|)$. Nierówność trójkąta daje żądane oszacowanie kombinacji liniowej. Jeżeli $|f|,|g|\le M$ na $A$ dla pewnego $M\ge1$, korzystamy z nierówności
\[
 |f(x)g(x)-f(y)g(y)|\le M|f(x)-f(y)|+M|g(x)-g(y)|.
\]
Wspólny promień dla dokładności $\varepsilon/(2M)$ wystarcza.
(iv) Z tożsamości różnicy odwrotności mamy
$|1/g(x)-1/g(y)|\le |g(x)-g(y)|/m^2$. Stosujemy jednostajną ciągłość $g$ dla dokładności $m^2\varepsilon$.
\end{proof}
Założenia ograniczoności w (iii) nie można ogólnie usunąć: iloczyn dwóch identyczności na $\R$ jest funkcją $x^2$. W (iv) sama niezerowość nie wystarcza, co pokazuje identyczność na $(0,1]$ i jej odwrotność z przykładu \ref{ma1:ex:reciprocal-not-uniform}.

\section{Ciągi Cauchy'ego i przedłużanie}
\begin{twierdzenie}
\label{ma1:thm:uniform-preserves-cauchy}
Jeśli $f:A\to Y$ jest jednostajnie ciągłe, gdzie $A\subseteq X$ i $X,Y$ są przestrzeniami metrycznymi, to obraz każdego ciągu Cauchy'ego w $A$ jest ciągiem Cauchy'ego w $Y$.
\end{twierdzenie}
\begin{proof}
Dla $\varepsilon>0$ wybieramy $\delta>0$ z jednostajnej ciągłości. Jeżeli $(x_n)$ jest ciągiem Cauchy'ego, istnieje $N$ takie, że $d_X(x_n,x_m)<\delta$ dla $n,m\ge N$. Wówczas $d_Y(f(x_n),f(x_m))<\varepsilon$ dla wszystkich takich indeksów.
\end{proof}
Odwrotność jest fałszywa: funkcja $x^2$ na $\R$ zachowuje ciągi Cauchy'ego, ponieważ każdy z nich jest zbieżny w $\R$, a obraz przez funkcję ciągłą jest zbieżny, zatem jest ciągiem Cauchy'ego. Nie jest ona jednak jednostajnie ciągła.

\begin{definicja}
\label{ma1:def:metric-closure}
Niech $X$ będzie przestrzenią metryczną i $A\subseteq X$. Domknięciem $A$ w $X$ nazywamy zbiór
$\overline A=\{p\in X:\forall r>0,\ B_X(p,r)\cap A\ne\varnothing\}$.
\end{definicja}
Każdy punkt $A$ należy do $\overline A$. Punkt $p$ należy do $\overline A$ wtedy i tylko wtedy, gdy jest granicą pewnego ciągu w $A$: wybieramy $x_n\in B_X(p,1/n)\cap A$ lub, w drugą stronę, korzystamy z definicji zbieżności. Ponadto $\overline A$ jest domknięty: dla $p\notin\overline A$ istnieje $r>0$ z $B_X(p,r)\cap A=\varnothing$. Jeżeli $q\in B_X(p,r/2)$, to $B_X(q,r/2)\subseteq B_X(p,r)$, zatem $q\notin\overline A$.

\begin{twierdzenie}[Przedłużanie jednostajnie ciągłych odwzorowań]
\label{ma1:thm:uniform-extension}
Niech $X$ będzie przestrzenią metryczną, $Y$ zupełną przestrzenią metryczną, $A\subseteq X$ oraz $f:A\to Y$ jednostajnie ciągłym odwzorowaniem. Istnieje dokładnie jedno ciągłe odwzorowanie $F:\overline A\to Y$ spełniające $F|_A=f$. Ponadto $F$ jest jednostajnie ciągłe. Jeżeli $f$ spełnia warunek H\"oldera z wykładnikiem $\alpha$ i stałą $C$, to $F$ spełnia ten warunek z tymi samymi $\alpha,C$.
\end{twierdzenie}
\begin{proof}
Dla $A=\varnothing$ także $\overline A=\varnothing$ i teza jest natychmiastowa. Niech $p\in\overline A$. Wybieramy ciąg $(x_n)$ w $A$ zbieżny do $p$. Jest on ciągiem Cauchy'ego, więc z twierdzenia \ref{ma1:thm:uniform-preserves-cauchy} ciąg $(f(x_n))$ jest ciągiem Cauchy'ego w $Y$. Zupełność $Y$ daje jego granicę; definiujemy $F(p)=\lim_n f(x_n)$.

Definicja nie zależy od wyboru ciągu. Jeżeli $y_n\to p$ w $A$, to $d_X(x_n,y_n)\le d_X(x_n,p)+d_X(p,y_n)\to0$. Z kryterium \ref{ma1:thm:uniform-sequential} odległości obrazów dążą do zera. Z nierówności trójkąta odległość ich granic jest nie większa niż suma trzech odległości: od pierwszej granicy do $f(x_n)$, od $f(x_n)$ do $f(y_n)$ i od $f(y_n)$ do drugiej granicy. Wszystkie trzy dążą do zera, więc granice są równe. Dla $p\in A$ można wybrać ciąg stały $x_n=p$, co daje $F(p)=f(p)$.

Pokażemy jednostajną ciągłość $F$. Dla $\varepsilon>0$ wybieramy $\delta>0$ z jednostajnej ciągłości $f$ dla dokładności $\varepsilon/2$. Niech $p,q\in\overline A$ i $d_X(p,q)<\delta/3$. Wybieramy $x_n,y_n\in A$ z $x_n\to p$, $y_n\to q$. Od pewnego indeksu $d_X(x_n,p)<\delta/3$ i $d_X(y_n,q)<\delta/3$, więc $d_X(x_n,y_n)<\delta$ oraz $d_Y(f(x_n),f(y_n))<\varepsilon/2$. Przechodząc do granicy w ostatniej nierówności, uzyskujemy $d_Y(F(p),F(q))\le\varepsilon/2<\varepsilon$. Przejście to wynika z oszacowania
\[
 |d_Y(f(x_n),f(y_n))-d_Y(F(p),F(q))|
 \le d_Y(f(x_n),F(p))+d_Y(f(y_n),F(q)).
\]
Jeżeli istnieje inne ciągłe przedłużenie $G$, to dla każdego $p\in\overline A$ wybór $x_n\to p$ w $A$ daje $G(p)=\lim_n G(x_n)=\lim_n f(x_n)=F(p)$. To dowodzi jednoznaczności wśród wszystkich ciągłych przedłużeń.

Na koniec, przy warunku H\"oldera, dla powyższych ciągów mamy $d_Y(f(x_n),f(y_n))\le C d_X(x_n,y_n)^\alpha$. Odległości argumentów dążą do $d_X(p,q)$ na mocy analogicznego oszacowania różnicy odległości. Funkcja potęgowa $t\mapsto t^\alpha$ jest ciągła na $[0,+\infty)$ dla $\alpha>0$; dla wykładników rzeczywistych uzasadnienie tej własności podamy przy konstrukcji funkcji elementarnych. Można tu również ominąć tę zależność: dla dowolnego $s>d_X(p,q)$ od pewnego indeksu $d_X(x_n,y_n)<s$, więc $d_Y(F(p),F(q))\le Cs^\alpha$. Kres dolny liczb $s^\alpha$ dla $s>d_X(p,q)$ wynosi $d_X(p,q)^\alpha$. Otrzymujemy żądane oszacowanie.
\end{proof}
\begin{przyklad}
\label{ma1:ex:uniform-extension-incomplete}
Nie można ogólnie pominąć zupełności $Y$. Niech $A=(0,1)$ w $\R$ i $Y=(0,1)$ z metryką zwykłą. Identyczność $f:A\to Y$ spełnia warunek Lipschitza ze stałą $1$, lecz nie ma ciągłego przedłużenia na $[0,1]$ o wartościach w $Y$: dla $x_n=1/(n+1)$ ciąg wartości dąży w $\R$ do $0$, a ciągłość przedłużenia wymagałaby granicy w $Y$.
\end{przyklad}
\begin{przyklad}
\label{ma1:ex:dense-extension}
Niech $A=\Q$, $X=Y=\R$ i $f:\Q\to\R$ będzie odwzorowaniem spełniającym warunek Lipschitza ze stałą $C$. Z gęstości $\Q$ w $\R$ mamy $\overline{\Q}=\R$. Twierdzenie \ref{ma1:thm:uniform-extension} daje dokładnie jedno ciągłe przedłużenie na $\R$; spełnia ono ten sam warunek Lipschitza. Wartość w liczbie niewymiernej $p$ otrzymujemy jako granicę $f(q_n)$ dla dowolnego ciągu liczb wymiernych $q_n\to p$.
\end{przyklad}

W następnym rozdziale rozważymy funkcje monotoniczne, których zachowanie graniczne wynika z porządku na prostej rzeczywistej.
